\section{Two-factor completion route: regular heads and transport of four long boundaries}
This analytic argument and the COMPLETE finite entrance below establish
all remaining two-factor minimal crossings. All parameters outside that
explicit finite entrance are covered by unbounded written proofs.
General C, A2 and original coefficient UI remain OPEN.

\paragraph{A valid window-length transport with a prefix-ratio premise.}
Let f be nonnegative UI, positive through L, strictly increasing to
its first mode M. Suppose it is LC at all centres K..M, K>=2,
and EVERY earlier positive ratio is at least f_(K+1)/f_K.
For A>=K set W_A(k)=sum(f_(k)..f_(k+A-1)). If W_K is pre-mode LC,
then every W_A is pre-mode LC. The same applies to actual signed
tails after positive clipping: any strictly growing positive window
triple has positive entering endpoint, so all its reads agree with f.

Proof: put m=k+A and D(r,s)=f_r*f_(s-1)-f_(r-1)*f_s. The exact minor is
\[
 H_A(k)=\sum_{s=k+2}^{m+1}D(k+1,s)
             +\sum_{r=k+2}^{m}D(r,m+1).
\]
If m<M, all terms D(r,m+1) are nonnegative by prefix compatibility
and tail LC. Also
\[
 D(r,m+2)\ge(f_{m+1}/f_m)D(r,m+1)\ge D(r,m+1).
\]
The newly added first and last terms are nonnegative. Thus H_A(k)
is nondecreasing in A while its right endpoint stays before M.
If m>=M, the window is concave because its entering increment is
nonpositive and outgoing increment nonnegative; its minor is automatic.
First modes of these forward windows do not increase with A: a strict
new next growth f_(k+A+1)>f_(k+1) implies the same old growth for
every smaller A, by strict input rise or the descending input tail.
Hence every required old minor lies in the seed's pre-mode band.
This proves the transport, with every premise retained.

\paragraph{All four long-capacity boundaries, analytically.}
For binary count u>=14 and single capacity B=u-c, c=-1,0,1,2,
put K=c+3. Its actual signed tail f starts at physical B.
The registered one-jump analytic theorem supplies LC at all centres
k>=K before its first mode (physical centres >=u+3).
The actual strict-rise theorem applies; its only possible early flat
would be at B+1, but f1>f0 here.

Put x=Fib(B), y=Fib(B+1), f0=x, f1=y-1. Until rank K-1, f is the
Fibonacci recurrence on these initial values. Its first two losses
from that same recurrence are a,b at K,K+1:
(a,b)=(2,2u+3) if alpha1, (1,u+3) if alpha2, (1,u+2) otherwise.
The reference first minor C=(y-1)^2-x*(x+y-1) is negative and
|C|<=7x/3, using y<=5x/3 and Cassini.
For i<K the exact prefix comparison is
\[
 f_i f_K-f_{i-1}f_{K+1}
   =(-1)^{i-1}Fib(K+1-i)C+b f_{i-1}-a f_i.
\]
Since reference ratios are <=2, b-2a>=u, f_(i-1)>=x, and
Fib(K)<=5, this is >=(u-35/3)x>0. At i=K use already proved tail LC.
All prefix ratios therefore dominate the first tail ratio.

For c=0,1,2, the seed window length K gives CI degree exactly u+1,
already registered for EVERY padding. Window transport closes ALL
larger capacities A>=K, including every remaining minimal crossing.
For c=-1, K2, the seed first minor instead has an explicit payment.
Writing eps=(-1)^(u+1), its first four values are
(x,y-1,x+y-a-1,x+2y-b-2), where x=Fib(u+1),y=Fib(u+2).
The first length-two window minor equals
\[
 eps+(b-a+1)x+(b-3a-2)y+a^2+3a+1-b\ge0.
\]
For alpha>=3, alpha2, alpha1 this becomes respectively
eps+(u+2)x+(u-3)y+3-u,
eps+(u+3)x+(u-2)y+2-u,
eps+(2u+2)x+(2u-5)y+8-2u, all positive at u>=14.
Every later window minor agrees with the window of the truncated
input f[1:], which is pre-mode LC; the signed-uniform theorem pays it.
Thus ALL maxcap>=u-2 minimal two-factor crossings are closed.

\paragraph{The regular-capacity cone u>=24, maxcap<=u-3.}
Write V=(A,B), A<=B<=u-3, h=u+delta, 3<=delta<=A-1.
Then delta<=u-8. The primitive spectral norm is <6/5, below the
registered 5/4 interior budget at ALL u>=16. Since B<=u-3, EVERY
natural head and last-junction centre is in that interior band.
For oldCount1, the last-junction scalar is (u^2-15u+8)/2>0,
giving gap3u/2, while its atom error is <=6u/5+2.
Thus the natural quotient and later positive primitive are LC.

Set m=floor(u/4)>=6 and truncate the original primitive below
b=min(m,delta). New head centres m+1..delta are paid as follows.
For s<=u/2, the ULC kernel ratio is >=Q_s/(4(m+2)^2), and
Q_s>=binom(u-1,m+1). Its atom coefficient is at most
(6/5)phi^(delta-m) <=(6/5)phi^(3m-5)
 <=(1/8)(17/4)^m.
For ALLm>=6,
\[
 binom(4m-1,m+1)>(1/2)(m+2)^2(17/4)^m.
\]
Base m6 is 245157>188575. This left binomial grows by >=8,
whereas the right side grows by <6. Thus every left-half rank-one
minor and row-zero read is paid.

For s>=u/2 put k=u-1-s>=7. Then its atom coefficient is at most
(6/5)phi^(k-6). The elementary bound
(2q)^k>=k(k-1)/11 for all k>=7 follows from exact bases7..10
and the ratio (2q)(k-1)>=k+1 at k>=10.
Consequently
\[
 Q_s/phi^k >= binom(u-1,k)/phi^k >=(u-1)(u-2)/44.
\]
Use binom(n,k)=n(n-1)binom(n-2,k-2)/(k(k-1)) and
binom(n-2,k-2)>=2^(k-2) in this half.
At ALLu>=24 this pays the ULC rank-one budget because
\[
 (u-1)(u-2)/44 >= (6/(5phi^6))
       (u+2)^2(u+q)/(4(u+1)).
\]
Row zero follows from the same bound; Q reads are strictly positive.
At the left read n=m, the adjacent binomial ratio loses at most a
factor three, still leaving (4/3)(m+2)^2 times the atom bound.
On the right read band put k=u-1-n>=6; binom(u-1,k)>=2^k while
its read atom is at most (6/5)phi^(k-7), so positivity is strict.
Thus the original positive primitive from b is globally LC.

\paragraph{A quaternary source cone.}
For U2=[z^j](1+z)^u/(1-z)^2, define its golden curvature Q2.
At u=4j, formal Lagrange with w=t(1+w)^4 gives diagonal
\[
 G(w),\quad G(z)=(1-phi*z)^2/((1-z^2)^2(1+z)(1-3z)).
\]
G1=2-2phi<0, G2=10-3phi>0; ALL Gn,n>=2 are nonnegative.
For the base h, hn=[5(2-phi)3^(n-1)+(2+3phi)(-1)^n]/4,n>=2.
Even terms are positive; each odd last block
(r+1)h1+r*h3=r(24-14phi)+2-2phi>=0 at r>=1.
The sharper coefficientwise estimate
w-t-4t^2 <=(11/4)(w^2-t^2)
follows from the nonnegative polynomial
(11w-4)(1+w)^8+4(1+w)^4+5w (its first nonzero degree is three).
Indeed four times the difference of the desired two sides is t^2 P(w)/w;
P(w)/w is a polynomial with nonnegative coefficients, so this is
a valid coefficientwise argument despite the rational substitution.
It pairs the only negative atom against G2; the residual coefficient
is (31-17phi)/2>=27/16. Since [t^k]w^2>=8 for k>=3,
Q2(4j,j)>2 for ALLj>=2. Extend by Pascal induction to ALLu>=4j;
do not discard negative coefficients after multiplication by (1+w)^e.
This pays the actual first capacity cuts and startup through b<=u/4.
The primitive cut covers every later source centre, retaining rank0
and the previous actual unit. Thus the regular cone is closed.

\paragraph{Exact remaining finite entrance.}
Old u<=13 is registered. Four long boundaries are analytic.
The regular argument above covers ALLu>=24. Hence only the complete
two-factor minimal profiles at u14..23 with maxcap<=u-3 are covered by
the following complete certificate, with EVERY padding through the analytic
clamp plus the infinite representative. There are exactly 245 profiles,
5860 complete padding hosts, 76201 integer comparisons and 43336
distinct integer margins. Independent positive suffix-box sums rebuild
all profiles, hosts, modes and minors. Ordinary Lean reduction checks
all 43336 margins in 170 chunks, with only propext. Finite QA elsewhere
is not a theorem premise.

\paragraph{Universal late-derivative conclusion and its limits.}
For EVERY sorted word $1^\alpha2^uV$ with $\alpha\ge1$,
nonempty CI capacities $V\ge3$, and $h=\sum(t-1)>u$, let
$W=\sum t$. Its ACTUAL positive derivative tail from physical $W$
is LC at every internal centre strictly before its FIRST tail mode.
Equivalently the actual derivative is LC at every pre-global-mode
centre $N\ge W+1$. Counts $u\le13$ use the old theorem. Every
other CI contains a deletion-minimal degree crossing; one-factor
long seeds and degree-$u+1,u+2$ seeds are old, all >=3-factor
crossings are registered, and the two-factor proof above closes
the last entrance. The exact arbitrary-insertion signed-tail identity
propagates through all remaining CI factors. At $u=0$ the earlier
registered signed-tail argument is immediate.

This theorem does not assert derivative LC at $h\le u$, where
an actual counterexample is already registered. It also does NOT
by itself prove the required earlier ordinary/cumulative C entrance,
A2 or original unamplified coefficient unimodality. Those remain OPEN.

\paragraph{Formal and primary-source scope.}
The standard formal Lagrange formula is Gessel, Theorem2.1.1
Eq.(2.1.4), \texttt{https://arxiv.org/html/1609.05988},
DOI\texttt{10.1016/j.jcta.2016.06.018}, already primary-verified.
Twenty-five new scoped unbounded algebraic, prefix/window, source,
norm, binomial-growth and half-band implications are kernel checked
in \texttt{TwoFactorCompleteEntrance.lean}. One further report
checks the full integer entrance in \texttt{TwoFactorEntranceCertificate.lean}.
Power-series substitution, binomial induction, window sum/mode
assembly, actual-word identities and insertion remain written proofs;
no end-to-end Lean composition theorem is claimed.
